\section{Replacement and integrated forcing}\label{sec:forcing}

We begin the local concentration proof by controlling how many
variables are forced along the clause process. The estimate is
integrated over prefix lengths: erasing one variable turns forcing into
a difference of satisfiability probabilities, and those differences
can be summed by telescoping.

\begin{samepage}
For the fixed \(k\ge3\), write
\[
 H_n=\inf\{m\ge1:F_{n,m}\notin\sat\}.
\]
The value \(+\infty\) is allowed in this definition, although
\(P_n(m)\le2^n(1-2^{-k})^m\) implies that \(H_n\) is finite almost
surely. Put
\begin{equation}\label{eq:statistic}
 L_k=2^{k+1},\qquad M_n=L_kn,\qquad
 V_n=\min\{H_n-1,M_n\},\qquad \mu_n=\frac{\EE V_n}{n}.
\end{equation}
Thus \(V_n\) is the last satisfiable index up to the cap \(M_n\), and
\begin{equation}\label{eq:survival}
 P_n(m)=\PP(V_n\ge m)\qquad(0\le m\le M_n).
\end{equation}
This identity includes \(m=M_n\). Later, when a clause is omitted, the
same endpoint convention will count every additional satisfiable prefix.
\end{samepage}

\subsection{Killing a solution set with one clause}

A variable of a satisfiable formula \(H\) is \emph{forced} if all
satisfying assignments of \(H\) give it the same value. Let \(b(H)\) be
the number of forced variables, and set \(b(H)=0\) if \(H\) is
unsatisfiable. This definition uses all solutions of the entire
formula, even when the formula contains clauses of different lengths.

For a formula \(H\) on \(v\) variables and \(1\le r\le v\), define
\begin{equation}\label{eq:kill}
 q_r(H)=\frac{\fall{b(H)}r}{2^r\fall vr},
 \qquad
 \fall ur=u(u-1)\cdots(u-r+1),
\end{equation}
where \(\fall ur=0\) for nonnegative integers \(u<r\).
Here a random \(r\)-clause uses \(r\) distinct uniform variables and
independent fair signs.

If \(H\) is satisfiable, \(q_r(H)\) is exactly the probability that
adding one independent random \(r\)-clause makes it unsatisfiable.
Indeed, every solution of \(H\) must falsify every literal in such a
clause. Each chosen variable must therefore be forced, and its sign
must oppose the forced value. Conversely those conditions make the
clause false in every solution. For unsatisfiable \(H\), the convention
\(q_r(H)=0\) records that there is no transition to count.

\subsection{Replacing a shorter clause}

The next estimate is uniform in the background formula. This
uniformity will allow successive replacements after conditioning on all
other clauses.

\begin{samepage}
\begin{lemma}[Replacing a \((k-1)\)-clause]\label{lem:replacement}
Let \(k\ge3\) and \(n\ge k+1\) be integers, and put
\[
 v=n-1,\qquad \eps=n^{-(k-1)/k},\qquad
 g=\lceil2kn^{1/k}\rceil.
\]
For any formula \(H\) on these \(v\) variables, let \(Q\) be an
independent random \((k-1)\)-clause and let \(R_1,\ldots,R_g\) be
independent random \(k\)-clauses, also independent of \(H,Q\). Then
\[
 \PP(H\wedge R_1\wedge\cdots\wedge R_g\notin\sat)
 \ge \PP(H\wedge Q\notin\sat)-\eps.
\]
\end{lemma}
\end{samepage}

The power comparison below is the Boolean forced-variable instance of
the codimension comparison in Carenini~\cite[Lemma~6.1]{Carenini}. The
proof records the block size and error needed for the present summation.

\begin{proof}
If \(H\) is unsatisfiable, both unsatisfiability probabilities are one.
Suppose it is satisfiable, and set \(x=q_{k-1}(H)\) and \(b=b(H)\).
When \(x\le\eps\), the right side after subtracting \(\eps\) is
nonpositive. Assume \(x>\eps\).

First \(b\ge k\). Otherwise
\[
 x\le\frac{1}{2^{k-1}\binom{n-1}{k-1}}
 \le\frac1n\le n^{-(k-1)/k}=\eps.
\]
The middle inequality follows from
\(\binom{n-1}{k-1}\ge n-1\), valid for \(1\le k-1\le n-2\), and
\(2^{k-1}(n-1)\ge n\). For \(b\ge k\), the exact killing probabilities
give
\[
 q_k(H)
 =x\frac{b-k+1}{2(v-k+1)}
 \ge \frac{x}{k}\frac{b}{2v}
 \ge \frac1k x^{k/(k-1)}.
\]
Here \(b-k+1\ge b/k\) and \(v-k+1\le v\). For the last inequality,
each factor \((b-j)/(v-j)\) is at most \(b/v\), so
\(x\le(b/(2v))^{k-1}\). Consequently,
\[
 gq_k(H)\ge2n^{1/k}x^{k/(k-1)}\ge2x,
\]
because \(x>n^{-(k-1)/k}\).

The block makes \(H\) unsatisfiable whenever at least one of its
clauses does so individually. These individual events are independent
conditional on \(H\). Hence
\[
 \PP(H\wedge R_1\wedge\cdots\wedge R_g\notin\sat)
 \ge1-(1-q_k(H))^g
 \ge1-e^{-2x}
 \ge x.
\]
For the last inequality, \(0\le x\le2^{-(k-1)}\le1/4\), and
\(1-e^{-2x}-x\) vanishes at zero and has positive derivative on that
interval.
\end{proof}

\subsection{Erasing one variable and summing over time}

The following estimate is the input to the omission argument in
Section~\ref{sec:concentration}. Unsatisfiable formulas contribute zero,
and the sum includes the empty prefix and the cap.

\begin{proposition}[Integrated forcing]\label{prop:forcing}
For every fixed integer \(k\ge3\),
\begin{equation}\label{eq:integrated}
 \sum_{j=0}^{M_n}\EE\frac{b(F_{n,j})}{n}=O_k(n^{1/k}),
 \qquad M_n=2^{k+1}n.
\end{equation}
\end{proposition}

\begin{proof}
It is enough to consider \(n\ge k+1\). We first prove an explicit
bound for every integer cap \(M\ge0\), and then take \(M=M_n\).
The case \(M=0\) is immediate, so assume \(M\ge1\).
By symmetry, the \(j\)-th
summand is the probability that \(F_{n,j}\) is satisfiable and \(x_n\)
is forced. Delete the clauses containing \(x_n\), and let \(d\) be
their number. Then \(d\sim\Bin(j,k/n)\). Conditional on their positions,
the untouched formula \(B\) is an independent proper \(k\)-SAT formula
on \(v=n-1\) variables with \(j-d\) clauses. Erasing the literal on
\(x_n\) from each deleted clause leaves independent uniform
\((k-1)\)-clauses \(Q_1,\ldots,Q_d\) on those \(v\) variables,
independent of \(B\).

If \(F_{n,j}\) is satisfiable and \(x_n\) is forced, then
\begin{equation}\label{eq:forced-reduction}
 B\in\sat,\qquad B\wedge Q_1\wedge\cdots\wedge Q_d\notin\sat.
\end{equation}
Indeed, any solution of the latter formula would extend to a solution
of \(F_{n,j}\) with either value of \(x_n\).

Keep the incidence positions fixed. Replace each \(Q_\ell\) in turn by
\(g\) fresh random \(k\)-clauses, with \(g,\eps\) as in
Lemma~\ref{lem:replacement}. At a replacement step, condition on \(B\),
all other remaining shortened clauses, and all previously inserted
blocks. The clause being replaced is still uniform and independent of
this background, so the lemma bounds the loss in unsatisfiability
probability by \(\eps\). Averaging at each step yields
\[
 \PP(B\wedge Q_1\wedge\cdots\wedge Q_d\notin\sat
       \mid\text{incidence positions})
 \le1-P_v(j-d+dg)+d\eps.
\]
The formula after all replacements has \(j-d+dg\) independent proper
\(k\)-clauses on \(v\) variables. Since \(B\notin\sat\) is contained in
the event on the left and has conditional probability \(1-P_v(j-d)\),
the conditional probability of \eqref{eq:forced-reduction} is at most
\begin{equation}\label{eq:replacement-difference}
 \Delta_{j,d}+d\eps,\qquad
 \Delta_{j,d}=P_v(j-d)-P_v(j-d+dg)\ge0.
\end{equation}

To sum this estimate, let
\(p_{j,d}=\PP(\Bin(j,k/n)=d)\). For \(0\le d\le j\le M\),
\[
 p_{j,d}\le\binom jd(k/n)^d\le\frac{(kM/n)^d}{d!}.
\]
For fixed \(d\), monotonicity of \(P_v\) and telescoping give
\begin{equation}\label{eq:telescoping}
 \sum_{j=d}^{M}\Delta_{j,d}
 =\sum_{t=0}^{M-d}\bigl[P_v(t)-P_v(t+dg)\bigr]\le dg.
\end{equation}
For example, after shifting the second sum by \(dg\), only the first
\(dg\) nonnegative terms can remain; if the shift exceeds the number
of terms, that number is already at most \(dg\).

Use the envelope \((kM/n)^d/d!\) only for the nonnegative
\(\Delta_{j,d}\) term. The error \(d\eps\) is averaged with its original
binomial weight, whose mean is \(kj/n\). Therefore
\begin{align}
 \sum_{j=0}^{M}\EE\frac{b(F_{n,j})}{n}
 &\le
 \sum_{d=0}^{M}\frac{(kM/n)^d}{d!}
       \sum_{j=d}^{M}\Delta_{j,d}
       +\eps\sum_{j=0}^{M}\frac{kj}{n}\notag\\
 &\le
 g\sum_{d=0}^{\infty}\frac{d(kM/n)^d}{d!}
       +\frac{kM(M+1)}{2n}\eps\notag\\
 &=\frac{kM}{n}e^{kM/n}g+\frac{kM(M+1)}{2n}\eps.
 \label{eq:forcing-general-cap}
\end{align}
At \(M=M_n=L_kn\), the exponential factor depends only on \(k\).
The last expression is therefore \(O_k(g+n\eps)=O_k(n^{1/k})\),
because \(g=O_k(n^{1/k})\) and \(n\eps=n^{1/k}\).
\end{proof}
